Stratification is an operation that refines a model to capture finer detail
about subpopulations.  For instance, we can stratify the SIR model by splitting the susceptible state $(S, \{(vac_T), (vac_F)\})$ into two subpopulations $(S, \{(vac_T)\})$ and $(S, \{(vac_F)\})$.  State refinements imply transition and parameter refinements.  For example, stratifying $S$ implies that the $inf$ transition becomes stratified to represent how two different subpopulations (vaccinated or not) become infected.  Modifying the $\beta$ parameter for the stratified $inf$ transitions reflects how vaccinated individuals become infected at a lower rate.

\begin{definition}
	A stratification operator

	\begin{eqnarray*}
		St&:& \Omega \rightarrow \Omega'\\
		St(\Omega, \Theta, (x, g)) &=& (\Omega', \Theta')
	\end{eqnarray*}

	defines a mapping from state variable $(x, g)$ to a set of state variables
	$\{(x, g_1),\ldots, (x, g_N)\}$.  We assume that there is a unique $v\in V$ where $({\cal X}(v),{\cal D}(v)) = (x, g)$. The result of applying the operator  is a new Petrinet
	$\Omega'$ and semantics $\Theta'$, as follows:

	\begin{itemize}
		\item $V_X' = V_X \backslash \{v\}\cup \{v_1, \ldots, v_N\}$ where $({\cal X}(v),{\cal D}(v)) = (x, g)$, $({\cal X}'(v_i),{\cal D}'(v_i)) = (x, g_i)$ for $i \in [1, N]$, $g = g_1 \cup \ldots \cup g_N$, and $g_i \cap g_j = \emptyset$ for $i\not=j$. $V_{Z'} = V_Z \backslash V_z(v) \cup \{v_{z'} |  z' \in St_z((x, g), \{(x, g_1),\ldots, (x, g_N)\})\}$ where $V_z(v) = \{z\in Z | (v_z, v) \in E \} \cup \{z\in Z | (v, v_z) \in E\}$ and $St_z((x, g), \{g_1,\ldots, g_N\})$ is the operator for stratifying a transition with respect to $x$.
		\item $E'= E \backslash E_Z(v) \cup E_{Z'}(v)$ where $E_Z(v) = \{(v, v_z) \in E \} \cup \{(v_z, v) \in E \}$ and $E_{Z'}(v) = \{(v_i, v_{z_i})\}$

		      Let $E' = E \backslash \left(\bigcup_{z \in Z(x,S)} E(z)\right) \cup \bigcup_{z \in Z(x,S)} E(z)\bigcup_{i=1}^M E(z_i)$, where For each set of state transitions $st(z_i)$, $E(z_i) = E^{in}(z_i) \cup E^{out}(z_i)$, $E^{in}(z_i) =\{((w, R), z_i) | ((w, R), z_i, (y, T)) \in st(z_i)\}$, and  $E^{out}(z_i) =\{(z_i, (y, T)) | ((w, R), z_i, (y, T)) \in st(z_i)\}$.

		\item $P' = P \backslash P_{St}^- \cup P_{St}^+$, where $P_{St}^-$ and $P_{St}^+$ are the respective parameters removed or added because of the stratification, as defined below.
		\item $D' = D$
		\item $X' = X$


		\item $Z' = Z \backslash \{z \in Z| v_z \in V_z(v), {\cal Z}(v_z) = z\} \cup \{z_i |  z_i \in St_z((x, g), \{(x, g_1),\ldots, (x, g_N), v_z \in V_z(v), {\cal Z}(v_z) = z \}\}$ is the set of transitions from $Z$ that do not involve $(x, g)$ and the transitions added by stratifying transitions in $Z$ involving $(x, g)$. The stratification defines a set of new transitions $\{z_1,\ldots,z_M\}$ by i) stratifying each state transition in $st(z)$ into a set of state transitions, ii) taking the cross-product of the state transition sets, and iii) defining $M$ new state transition sets $st(z_i)$ each according to an element of the cross-product.  Stratifying a state transition $((w, g_R), z, (y, g_T)) \in st(z)$, according to the new states $\{(x, g_1),\ldots, (x, g_N)\}$ corresponds to one of the following cases:
		      \begin{enumerate}
			      \item If $w, y \not= x$, then the resulting set of state transitions is $\{((w, g_R), z, (y, g_T))\}$ when the state transition is unchanged because it does not depend on $x$.
			      \item If $w = x, y \not= x$, then the resulting set of state transitions is $\{((x, g_R\cap g_i), z, (y, g_T))| (x,g_i) \in \{(x, g_1),\ldots, (x, g_N)\}\}$ when the	source state is $x$.
			      \item If $w \not= x, y = x$, then the resulting set of state transitions is $\{((w, g_R), z, (x, g_T\cap g_i))| (x,g_i) \in \{(x, g_1),\ldots, (x, g_N)\}\}$ when the target state is $x$.
			      \item If $w,y = x$, then the resulting set of state transitions is $\{((x, g_R\cap g_i), z, (x, g_T\cap g_j))| (x,g_i), (x,g_j) \in \{(x, g_1),\ldots, (x, g_N)\}\}$ when the source and target state are $x$.
		      \end{enumerate}

		      In cases 3) and 4) above, the flow from each source state is divided
		      across $N$ state transitions.  The sum of the rates
		      associated with the transitions for each source state must equal the
		      original rate associated with $((w, g_R), z, (y, g_T))$.  In these cases, the new rates
		      preserve the flow by introducing a parameter $p_j$ for each transition and
		      ensuring that for 3) $\sum_{j=1}^N  p_{j} = 1$ and for 4) $\sum_{j=1}^N
			      p_{ij} = 1$ for each $i = 1\ldots N$.  These parameters represent the
		      probability of transitioning from $(w, g_R)$ to $ (x, g_T\cap g_i)$  or from $(x, g_R\cap g_i)$ to $(x, g_T\cap g_j)$, respectively.
		\item ${\cal P}'$
		\item ${\cal X}'$
		\item ${\cal D}'$
		\item ${\cal I}'$
		\item ${\cal Z}'$
		\item Let ${\cal R}'({\bf p}', {\bf x}', z') ={\cal R}({\bf p}, {\bf x}, z)[(x, S)/(x, g_i), \ldots, (x, S)/(x, g_N)]$  , for each $z \in Z'$ substitute occurrences of $(x, S)$ with $(x, g_i)$

		      \begin{eqnarray*}\label{eqn:stratified-rate}
			      {\cal R}'({\bf p}', {\bf x}', z') = \left\{
			      \begin{array}{lll}
				      {\cal R}({\bf p}, {\bf x}, z)     & : z \not\in Z(x, S)                \\
				      {\cal R}({\bf p}', (x, g_i), z_i) & : z \in Z(x, S), ((x, S), z) \in E
			      \end{array}
			      \right.
		      \end{eqnarray*}

	\end{itemize}
\end{definition}

For example, Example \ref{ex:base} is a stratification of the SIR model that
captures two susceptible groups that have different infection rates $\beta_1$
and $\beta_2$.  Stratifications also capture how the subpopulations change and
interact.  For example, if the subpopulations refer to whether individuals have
been vaccinated or not, then it is possible for unvaccinated individuals to
become vaccinated. Similarly, individuals becoming infected may be exposed to
susceptible individuals that have been vaccinated or not, each with a different
infection rate.  Transitions between strata can occur in conjunction with
another transition (e.g., simultaneously becoming infected and vaccinated), or
via ``self'' transitions (e.g., transitioning between susceptible groups when
vaccinated). Stratification can also be partial, meaning that states like $S_1$
and $S_2$ are stratified, but $I$ and $R$ are not.  This is useful when, for
example, the recovery rate is the same for all subpopulations and there is no
benefit to stratifying $I$ or $R$.


\begin{definition}
	A transition component $t \in T(z)$ is a pair $(x_i, x_o)$ where $(x_i, z)
		\in E$ and $(z, x_o) \in E$.
\end{definition}

The transition components express the possible one to one mappings of inputs and
outputs.  There may be several such mappings, but we assume for the sake of
simplicity, that there is a single unique mapping.  This single mapping includes
at most one ``natural'' transition component and zero or more ``noop''
components.  Natural transitions involve a change of state where the input and
output state are not equal.  The noop transitions require that the input and
output states are identical.  This assumption imposes a restriction on the
possible transitions needed to have a single set of transition components
describing a transition.


\begin{definition}
	Stratifying a transition component $t = (x_{in}, x_{out}) \in T(z)$ defines
	a set of transition component and strata transition pairs $\{((x_{in},
		x_{out}), (s_i(x_{in}), s_i(x_{out})) | {\cal D}= [D_0 | D_1| \ldots,
		V(s_i(x_{in})) = V(s_i(x_{in})) \cap_a D_i, V(s_i(x_{out})) =
		V(s_i(x_{out})) \cap_a D_i \}$
\end{definition}



%%% Local Variables:
%%% mode: LaTeX
%%% TeX-master: "main"
%%% End:
